\section{Conductor reduction and arbitrary contact multiplicities}
\label{sec:conductor-primary}
The quartic theorem determines an embedded structure whose thickness
grows with the symmetric degree. We next give a different global
classification, on non-hyperplane Grassmannian families of unbounded
codimension. Here the failure scheme stabilizes after a degree wall,
and the contact multiplicities determine every primary exponent.

\subsection{An exact reduction to a finite algebra}
Fix an effective divisor $Z$ of degree $d$ on $\Pj^1$, let $g$ be
its homogeneous equation, and suppose $n\ge2d-1$. Write
\[
 W=H^0(\mathcal O(n)(-Z))=gV_{n-d},\qquad
 R_j=H^0(Z,\mathcal O_Z(jn)).
\]
Let $B_{Z,s}$ be the open subset of $\Gr(s,R_1)$ consisting of
subspaces generating the invertible $\mathcal O_Z$-module
$\mathcal O_Z(n)$. For $A\in B_{Z,s}$ put
\[
 U_A=\operatorname{res}_Z^{-1}(A)\subset V_n,
 \qquad \dim U_A=n-d+1+s.
\]
This realizes $B_{Z,s}$ as the open base-point-free part of the
Grassmannian of subseries containing $W$.

\begin{theorem}[Conductor reduction, including all Fitting ideals]
\label{thm:conductor-reduction}
For $n\ge2d-1$, $1\le s\le d$, and $m\ge2$, restriction induces
a canonical isomorphism of coherent sheaves on $B_{Z,s}$,
\begin{equation}\label{eq:conductor-cokernel}
 \operatorname{coker}\big(\Sym^m\mathcal U\longrightarrow
                         V_{mn}\otimes\mathcal O\big)
 \simeq
 \operatorname{coker}\big(\Sym^m\mathcal A\longrightarrow
                         R_m\otimes\mathcal O\big).
\end{equation}
In particular all Fitting ideal sheaves, with the same Fitting
index, agree. These assertions commute with arbitrary base change.
They also hold relatively when $Z$ varies in a flat family of
effective degree-$d$ divisors. The reduction is effected by the
identity
\begin{equation}\label{eq:conductor-image}
 WU_A^{m-1}=gV_{mn-d}.
\end{equation}
\end{theorem}
\begin{proof}
Choose $a\in U_A$ whose restriction to $Z$ is nowhere zero. Such
an element exists: at each of the finitely many support points,
vanishing is one proper linear condition on $A$. Avoid their union.
The subspace $U_A$ is base-point-free on $\Pj^1$, since the common
zero scheme of $W$ is $Z$ and $a$ is a unit there.

We first prove $WU_A=gV_{2n-d}$. Products in $W^2$ span
$g^2V_{2n-2d}$. After division by $g$, restriction identifies
\[
 gV_{2n-d}/g^2V_{2n-2d}
 \simeq H^0(Z,\mathcal O_Z(2n-d)).
\]
The restriction $V_{n-d}\to H^0(Z,\mathcal O_Z(n-d))$ is onto:
its obstruction group is $H^1(\mathcal O(n-2d))=0$ under the
stated inequality. Multiplication by $a|_Z$ is an isomorphism of
these invertible modules. Thus $Wa$ surjects onto the displayed
quotient, while $W^2$ supplies its kernel. This proves the initial
identity, for nonreduced as well as reduced divisors.

By Lemma~\ref{lem:hyperplane-propagation}, whose proof applies to
any base-point-free subspace, $U_AV_e=V_{n+e}$ for $e\ge n-1$.
Starting with $e=2n-d$ and iterating proves
\eqref{eq:conductor-image} for every $m\ge2$.
The target restriction sequence
\[
 0\longrightarrow gV_{mn-d}\longrightarrow V_{mn}
       \longrightarrow R_m\longrightarrow0
\]
is exact. Its kernel is contained in the image of multiplication
by \eqref{eq:conductor-image}, and the image of multiplication
in the quotient is exactly $A^m$. Taking cokernels proves the
asserted isomorphism on fibres.

The same proof gives a sheaf isomorphism, rather than merely equality
of fibre dimensions. Locally on the parameter space choose a section
$a$ with invertible restriction. The restriction maps above are
surjections of vector bundles, and multiplication by $a|_Z$ is a
bundle isomorphism. The image of
$\mathcal W\otimes\Sym^{m-1}\mathcal U$ therefore contains the
whole target kernel as a sheaf, by the preceding surjections and
induction. Quotienting gives \eqref{eq:conductor-cokernel} directly.
For a varying divisor, the same cohomology vanishings make all the
restriction sequences locally free and compatible with base change.
Fitting ideals depend only on a presentation of the cokernel and
commute with base change; see \cite[Tags 07Z8, 07ZA]{StacksFitting}.
\end{proof}

\subsection{Pencils in an arbitrary contact algebra}
We now take $s=2$ and $d\ge3$. Trivialize the line bundle on $Z$
and write
\[
 Z=\sum_{i=1}^r e_i p_i,\quad p_i\ne p_j\ (i\ne j),\qquad
 R=H^0(Z,\mathcal O_Z)
       \simeq\prod_{i=1}^r\C[t_i]/(t_i^{e_i}).
\]
On the frame space of $B_{Z,2}$, choose a frame $(a,b)$ with
$a\in R^\times$. Dividing by $a$, write
\begin{equation}\label{eq:residue-pencil-coordinates}
 z=b/a,\qquad
 z_i=\lambda_i+c_{i1}t_i+\cdots+c_{i,e_i-1}t_i^{e_i-1}.
\end{equation}
The coordinates of $a$ and of $z$ form regular coordinates on an
open subset of the $2d$-dimensional frame space. The rank-two
condition excludes the locus where $z$ is a scalar in $R$.
These frame charts cover every base-point-free pencil.

\begin{theorem}[Primary classification for arbitrary contact pencils]
\label{thm:contact-pencil-primary}
Let $n\ge2d-1$, $d\ge3$, and let $D_{m,Z}$ be the maximal-minor
failure scheme of $\Sym^m\mathcal U\to V_{mn}\otimes\mathcal O$
on $B_{Z,2}$. Its subseries have codimension $d-2$ in $V_n$.
For $2\le m<d-1$ the ideal is zero, so $D_{m,Z}=B_{Z,2}$ as a
scheme. For every $m\ge d-1$, its ideal on the frame chart
\eqref{eq:residue-pencil-coordinates} is principal, generated, up
to a unit, by
\begin{equation}\label{eq:contact-index-form}
 \Delta_Z(z)=
 \prod_{i:e_i\ge2}c_{i1}^{\binom{e_i}{2}}
 \prod_{i<j}(\lambda_j-\lambda_i)^{e_i e_j}.
\end{equation}
In particular it is independent of $m$ in this range.

Let $\mathfrak R_i$ be the divisor $c_{i1}=0$ for $e_i\ge2$,
and let $\mathfrak C_{ij}$ be the divisor $\lambda_i=\lambda_j$.
They descend to distinct irreducible Cartier divisors on $B_{Z,2}$.
The complete primary decomposition is
\begin{equation}\label{eq:contact-primary-decomposition}
 \mathcal I_{D_{m,Z}}
 =\bigcap_{i:e_i\ge2}\mathcal I_{\mathfrak R_i}^{\binom{e_i}{2}}
   \ \cap\!\bigcap_{i<j}\mathcal I_{\mathfrak C_{ij}}^{e_i e_j}.
\end{equation}
Its associated primes are exactly those of these divisors; there
are no embedded associated primes. All components have codimension
one in $B_{Z,2}$.

At an arbitrary pencil, let $r_i$ be the least $j\ge1$ with
$c_{ij}\ne0$, and put $r_i=e_i$ if $z_i$ is constant. Define
\begin{equation}\label{eq:contact-minimal-polynomial-degree}
 q(z)=\sum_{\lambda\in\{\lambda_1,\ldots,\lambda_r\}}
          \max_{i:\lambda_i=\lambda}
                         \left\lceil\frac{e_i}{r_i}\right\rceil.
\end{equation}
Then for every $m\ge2$ the multiplication corank is
\begin{equation}\label{eq:contact-all-coranks}
 \operatorname{corank}\mu_{m,U_A}=d-\min\{m+1,q(z)\}.
\end{equation}
All statements hold over the ordered configuration space of the
$p_i$ with their fixed multiplicities, not only for a fixed divisor.
\end{theorem}

\begin{lemma}[Power-basis determinant]
\label{lem:confluent-power-basis}
In the algebra $R$ above, the determinant of the columns
$1,z,\ldots,z^{d-1}$, written in the bases
$1,t_i,\ldots,t_i^{e_i-1}$, is
\eqref{eq:contact-index-form}, up to a nonzero constant depending
only on the ordering of the bases. Over the polynomial coefficient
ring, every column $z^k$, $k\ge d$, is a linear combination of
these first $d$ columns.
\end{lemma}
\begin{proof}
Factor the evaluation map on polynomials $P$ of degree less than
$d$ into
\[
 P\longmapsto
    (P(\lambda_i+u_i)\bmod u_i^{e_i})_i
 \longmapsto
    (P(z_i(t_i))\bmod t_i^{e_i})_i.
\]
For the first map, use the monic Newton basis
\[
 (T-\lambda_i)^j\prod_{h<i}(T-\lambda_h)^{e_h},
       \qquad 0\le j<e_i,
\]
ordered by blocks $i$ and then $j$. The change from the monomial
basis is triangular with diagonal one. Evaluation on the jet
blocks is block triangular, with determinant
$\prod_{h<i}(\lambda_i-\lambda_h)^{e_h e_i}$.
For the second map, substitution $u_i=z_i(t_i)-\lambda_i$ is
triangular in powers of $t_i$, with diagonal
$1,c_{i1},c_{i1}^2,\ldots,c_{i1}^{e_i-1}$.
Multiplying these determinants proves the formula as a polynomial
identity, also when some differences or some $c_{i1}$ vanish.

Multiplication by $z$ is an endomorphism of the free rank-$d$
coefficient module $R$. Its monic characteristic polynomial has
coefficients in the same polynomial ring. Cayley--Hamilton, applied
to the element $1\in R$, expresses $z^d$ in the first $d$ powers;
iteration does the same for all higher powers. No division by the
determinant is involved.
\end{proof}

\begin{proof}[Proof of Theorem~\ref{thm:contact-pencil-primary}]
By Theorem~\ref{thm:conductor-reduction}, multiplication has the
same cokernel as the columns
$a^m,a^{m-1}b,\ldots,b^m$ in $R$. Multiplication by the unit
$a^{-m}$ turns these into $1,z,\ldots,z^m$. If $m+1<d$, this
matrix has no $d$-minors, so its zeroth Fitting ideal is zero.
If $m+1\ge d$, Cayley--Hamilton expresses every column after the
first $d$ in the earlier ones over the base ring. Multilinearity
therefore makes every maximal minor a multiple of their determinant,
which is itself one of the minors. Lemma~\ref{lem:confluent-power-basis}
proves the exact ideal formula.

The factors appearing in \eqref{eq:contact-index-form} are distinct
linear polynomials in regular frame coordinates. In the local
polynomial ring, their powers are primary and their product equals
their intersection, by unique factorization. This is an irredundant
primary decomposition, so there are no further associated primes.
These assertions descend from the frame bundle. More explicitly,
$\lambda_j-\lambda_i$ is a unit multiple of the determinant of
the two evaluation rows of the frame, and $c_{i1}$ is a unit
multiple of the determinant of its value and first-jet rows at
$p_i$. These define the Schubert divisors where a pencil fails to
map isomorphically to a fixed two-dimensional quotient of $R$.
Each is irreducible: it is the incidence divisor of pencils meeting
the fixed codimension-two kernel. The open base-point-free part is
nonempty for every factor when $d\ge3$. The different pairs and
jet quotients give different divisors. Frame changes multiply
local equations by units, and faithfully flat descent gives
\eqref{eq:contact-primary-decomposition} on $B_{Z,2}$.

For the corank calculation, the minimal polynomial of $z_i$ in
$\C[t_i]/(t_i^{e_i})$ is
$(T-\lambda_i)^{\lceil e_i/r_i\rceil}$. Indeed
$(z_i-\lambda_i)^k$ has order $kr_i$ until it vanishes.
The minimal polynomial of $z$ in the product algebra is the least
common multiple of these polynomials, whose degree is exactly
$q(z)$. The powers $1,z,\ldots,z^{q(z)-1}$ are linearly independent,
and all later powers lie in their span. This proves the rank and
corank formula. The proof is unchanged as the support points vary
in ordered distinct-support charts; changing local parameters
multiplies the determinant factors by units.
\end{proof}

\begin{corollary}[Complete local singularities and infinitesimal layers]
\label{cor:contact-local-structure}
For $m\ge d-1$, at a point of $D_{m,Z}$ let $\mathcal B$ be the partition of support indices
by equality of the values $\lambda_i$, and let
$J=\{i:e_i\ge2,\ c_{i1}=0\}$. On its smooth frame cover the
completed failure germ is a smooth factor times
\begin{equation}\label{eq:contact-weighted-arrangement}
 \prod_{i\in J}x_i^{\binom{e_i}{2}}
 \prod_{B\in\mathcal B}\ \prod_{i<j,\ i,j\in B}
               (y_j-y_i)^{e_i e_j}=0.
\end{equation}
Thus this equation determines the completed germ on $B_{Z,2}$
up to adding four smooth variables, the dimension of the frame
group. Its nilpotency index is the largest exponent of a factor
present in \eqref{eq:contact-weighted-arrangement}, with index one
for a reduced local ring.

More precisely write the local equation as $\Delta=\prod L_\alpha^{b_\alpha}$
and put $H=\prod L_\alpha$. For $k\ge1$, the successive powers
of the nilradical $N=(H)/(\Delta)$ satisfy
\begin{equation}\label{eq:contact-nilradical-layers}
 N^k/N^{k+1}
 \simeq R_{\rm loc}\big/
 \left(H,\prod_\alpha L_\alpha^{\max\{b_\alpha-k,0\}}\right),
\end{equation}
where the isomorphism sends $1$ to the class of $H^k$.
For reduced $Z$, the entire scheme is reduced; the collision
strata are indexed by set partitions with $q\ge2$ blocks, have
codimension $d-q$, and have multiplication corank $d-q$ when
$m\ge d-1$. A single pair collision is smooth, disjoint pair
collisions are ordinary normal crossings, and a block of size at
least three has the corresponding braid-arrangement germ.
\end{corollary}
\begin{proof}
All factors not vanishing at the chosen pencil are units. Translate
the remaining value coordinates and first derivatives to their
values at that pencil. The frame coordinates of $a$, the higher
jets, and blockwise common values are smooth parameters; this gives
\eqref{eq:contact-weighted-arrangement}. Since $H^k$ is divisible
by $\Delta$ exactly when $k\ge\max b_\alpha$, the claimed index
follows. Unique factorization gives
$(\Delta):H^k=(\prod L_\alpha^{\max\{b_\alpha-k,0\}})$,
which proves the layer formula. For reduced $Z$ all exponents are
one. Equalities of the values in a block of size $v$ impose $v-1$
independent equations. Their total is $d-q$; the free frame action
preserves this codimension. The local smoothness and crossing
assertions follow from the displayed product of differences.
\end{proof}

For example, for $Z=3p+2q$, $n\ge9$ and $m\ge4$, the subseries
have codimension three and the complete failure equation is
$c_{p1}^{3}c_{q1}(\lambda_q-\lambda_p)^6$. Its three primary
multiplicities are $3,1,6$; no embedded prime is suppressed by
passing to the residue calculation. For a reduced divisor with
arbitrarily many points, the same theorem gives all collision
components and all their intersections on a family of subseries
of arbitrarily large codimension. The finite-algebra determinant
and Cayley--Hamilton are classical; the content specific to
multiplication failure is the exact conductor identity and the
resulting identification of the full ideal sheaf, including these
primary structures and the degree wall.
